\newpage
\section*{\textsc{General Conclusion}}
\addcontentsline{toc}{section}{General Conclusion}
\rhead{\itshape General Conclusion}
\lhead{}
\cfoot{\thepage}
\pagestyle{fancy}
\renewcommand{\headrulewidth}{0.4pt}

\subsection*{Summary of Contributions}

This report presented \textbf{ExamGen}, an end-to-end bilingual exam generation
platform developed in response to a concrete problem in Algerian higher education:
instructors managing bilingual (French--Arabic), unstructured PDF archives with no
machine-readable indexing and no automated quality assurance for generated questions.

The thesis argued that a \textit{schema-enforced, RAG-augmented pipeline} combining
intelligent document processing, language-aware hybrid retrieval, and an independent
post-generation quality evaluator can produce examination questions that are
factually grounded, cognitively calibrated, and stylistically consistent with
historical practice --- without requiring model fine-tuning.
The experimental evaluation on a real corpus of 137 documents across five subjects
confirms this thesis on all four dimensions:

\begin{enumerate}
  \item \textbf{Bilingual IDP.}
  \texttt{docParser}'s five-stage pipeline (Docling, LLM extraction, grounding
  validation, Bloom classification, VLM diagram transcription) successfully processed
  90.5\% of documents, producing 86 validated Exam JSON files and 704 Bloom-labelled
  questions. The 24.1\% rejection rate is explainable: 21 answer-key sheets (no
  question text) and 12 low-resolution scans (below minimum OCR quality).

  \item \textbf{Language-aware hybrid RAG.}
  Routing Arabic Law queries to AraBERT and French queries to multilingual MiniLM,
  combined with BM25 and Reciprocal Rank Fusion, achieved 83.8\% overall retrieval
  effectiveness --- a 10-point improvement over the non-routed baseline on Arabic
  queries. This confirms the necessity of dedicated Arabic embeddings for low-resource
  Algerian legal content.

  \item \textbf{Schema-enforced generation.}
  The \texttt{instructor}-enforced structured output, combined with a dual-prompt
  architecture (style-anchor separation from RAG content), achieved 87.5\% factual
  accuracy across all five subjects --- consistent with the upper range reported by
  \citet{mucciaccia2025mcq} for prompt-engineered systems with retrieval augmentation.

  \item \textbf{Independent quality evaluation.}
  The 4-axis post-generation evaluator (factual grounding, Bloom alignment, quality
  verdict, difficulty calibration) confirmed 79.2\% cognitive-level accuracy,
  motivating the independent Stage-4 Bloom classifier rather than trusting the
  generator's self-reported level --- as recommended by \citet{meissner2024itemforge}.
\end{enumerate}

\subsection*{Lessons from Iterative Development}

The iterative design process was itself a methodological contribution.
Three concrete prototype failures --- absent structured-output enforcement, low
instruction-following fidelity on free-tier LLMs, and topic-blending from
dense-only retrieval --- each led directly to a documented architectural decision
(the \texttt{instructor} library, the DeepSeek-V4-Flash upgrade, and hybrid RRF
retrieval respectively). This failure-driven iteration mirrors the approach
advocated by \citet{isley2025quality}: continuous refinement guided by measurable
quality metrics.

\subsection*{Limitations}

Three limitations bound the scope of the current results:

\begin{enumerate}
  \item \textbf{Law corpus sparsity.}
  The Arabic Law corpus (11 JSON files, 66 questions) is the smallest in the dataset.
  This explains Law's lowest scores across all four evaluation axes.
  Expanding the corpus --- through additional lecturer contributions or data
  augmentation --- is the most direct path to closing this gap.

  \item \textbf{No human-labelled ground truth.}
  Factual accuracy and Bloom alignment are judged by the LLM evaluator, not by
  human subject-matter experts. This makes the 87.5\% factual accuracy figure an
  upper-bound estimate; \citet{isley2025quality}'s IRT-based validation against
  real student cohorts would provide a lower-bound ground truth.

  \item \textbf{API dependency.}
  The generation and evaluation stages rely on cloud LLM APIs
  (DeepSeek-V4-Flash, Gemini-2.0-Flash, GPT-OSS-120B).
  Connectivity or pricing changes could disrupt production use in a
  university environment with restricted internet access.
\end{enumerate}

\subsection*{Future Work}

Three directions would substantially extend this work:

\paragraph{(1) Unconditional OCR for low-resolution scans.}
The 12 rejected low-resolution PDFs could be recovered by activating Tesseract
or PaddleOCR unconditionally, at an estimated cost of 3--5~minutes per page on
CPU. A quality threshold on the OCR confidence score would gate whether the
result feeds into Stage~2.

\paragraph{(2) IRT-based psychometric validation.}
Deploying the generated exams in real examination sessions would enable
Item Response Theory calibration \cite{isley2025quality}: difficulty and
discrimination parameters estimated from student response vectors provide
an objective, human-grounded quality signal currently absent from the evaluation.

\paragraph{(3) On-device fine-tuned model.}
Fine-tuning a compact model (7B parameters, e.g.\ Qwen2.5-7B or Mistral-7B)
on the 704-question ExamGen dataset would eliminate API dependency and enable
offline operation --- a practical requirement for Algerian university networks.
The existing 4-axis evaluator provides ready-made preference signal for
DPO or RLHF alignment.
